% 问题三·子问题二：M1 的 N-D-Q_B 条件优化
% 本文件为正文片段，遵循 paper/main.tex 的唯一入口约定；不定义文档环境。
% 证据运行：q3-q-conditional-20260925-r01、q3-q-robustness-20260925-r01、
% q3-optimization-robustness-20260926-r01、q3-answer-package-20260926-r01。

\subsection{问题三子问题二：M1 的 $N$--$D$--$Q_B$ 条件优化}
\label{sec:q3-subproblem2}

\subsubsection{原生质量变量与问题分析}

本小问在M1模型的基础上增加质量评分决策变量。记为
$Q_B$；决策变量为模型规模 $N$、训练数据量 $D$ 和原生质量分数 $Q_B$。其中
$N,D$ 分别以十亿参数和十亿 token 计，B1 的观测支持域为
\begin{equation}
 N\in[0.070542,11.965825],\qquad
 D\in[0.134,299.893],
 \label{eq:q3-sp2-support}
\end{equation}
而 B6 原生质量分数的支持域为 $Q_B\in[0.1,1.0]$。

\subsubsection{M1 目标函数与质量成本}

冻结的 B6 M1 条件损失写为
\begin{equation}
 \widehat L_{\mathrm{M1}}(N,D,Q_B)
 =E+A N^{-\alpha}+B D^{-\beta}+G(1-Q_B),
 \label{eq:q3-sp2-m1-loss}
\end{equation}
其中本轮参数为
\begin{equation}
 \begin{aligned}
 E&=1.4892660413, & A&=0.5401729883, & B&=1.3167948616,\\
 G&=0.3622474350, & \alpha&=0.2790539239, & \beta&=0.2828029135.
 \end{aligned}
 \label{eq:q3-sp2-m1-parameters}
\end{equation}

给定上下文长度 $L_{\mathrm{ctx}}$，基础规模成本为
\begin{equation}
 C_{\mathrm{base}}=10^{18}ND(6+\eta L_{\mathrm{ctx}}),
 \qquad \eta=2\times10^{-4},
 \label{eq:q3-sp2-base-cost}
\end{equation}
相对于基准质量 $Q_0=0.1$ 的质量增量成本为
\begin{equation}
 C_Q=10^9D\,[g(Q_B)-g(Q_0)].
 \label{eq:q3-sp2-quality-cost}
\end{equation}
本文并列比较三类成本假设：
\begin{equation}
 \begin{aligned}
 g_{\mathrm{exp}}(Q)&=10^7\exp(6Q),\\
 g_{\mathrm{pow}}(Q)&=5\times10^9Q^4,\\
 g_{\mathrm{log}}(Q)&=2\times10^9\log(1+10Q).
 \end{aligned}
 \label{eq:q3-sp2-cost-families}
\end{equation}
三类 $g$ 是敏感性分析中的成本形状，不是由真实训练账单验证的成本定律。
在预算 $C$ 下，求解问题为
\begin{equation}
 \begin{aligned}
 \min_{N,D,Q_B}\quad &\widehat L_{\mathrm{M1}}(N,D,Q_B)\\
 \mathrm{s.t.}\quad &C_{\mathrm{base}}(N,D;L_{\mathrm{ctx}})+C_Q(D,Q_B)\le C,\\
 &N,D\text{ 满足式~\eqref{eq:q3-sp2-support}},\qquad Q_B\in[0.1,1.0].
 \end{aligned}
 \label{eq:q3-sp2-opt}
\end{equation}
质量成本占比记为
\begin{equation}
 \rho_Q=\frac{C_Q}{C_{\mathrm{base}}+C_Q}.
 \label{eq:q3-sp2-quality-share}
\end{equation}
$\rho_Q$ 只表示当前成本假设下的模型内比例，不是现实工程账单中的实测占比。

\subsubsection{求解方法}

数值求解采用 $(\log N,\log D,Q_B)$ 参数化。对预算约束先按预算值缩放，
再使用带显式不等式约束的 SLSQP，以避免 $10^{19}$--$10^{24}$ 量级的原始
约束造成有限差分的尺度失衡。每个情景使用六个确定性起点，覆盖低值、几何
中点、高值以及不同的 $Q_B$ 初始值；保留优化器收敛、预算可行且目标值最小的
候选，并记录所有起点的成功状态、目标值和边界标记。

当前问题的决策维度低、目标函数光滑、约束边界明确，解析结构与多起点局部优化更容易逐情景复核。求解后将状态分为
$Q_B=0.1$ 或 $Q_B=1.0$ 的端点、未触及边界的内点，以及触及 $N$、$D$ 或
$Q_B$ 支持边界的边界解。

\subsubsection{原生 $Q_B$ 情景结果}

本轮覆盖三类质量成本、五档上下文和三档预算，共
$3\times5\times3=45$ 个原生 $Q_B$ 情景。上下文和预算集合分别为
\[
 L_{\mathrm{ctx}}\in\{2048,4096,8192,32768,131072\},\qquad
 C\in\{10^{19},10^{22},10^{24}\}\ \mathrm{FLOPs}.
\]
每一行输出 $Q_B^*$、$N^*$、$D^*$、$L^*$、$\rho_Q$、预算是否活跃、
边界状态和多起点收敛信息；完整 45 行保存在
\texttt{experiments/runs/q3-q-conditional-20260925-r01/tables/}
下的情景表中。

按上下文取平均，低预算 $C=10^{19}$ 下，指数型、幂函数型和对数渐进型成本
分别给出 $Q_B^*=0.7243$、$0.6103$ 和 $1.0000$；预算提高到
$10^{22}$ 或 $10^{24}$ 后，三类成本在当前参数下均通常把 $Q_B^*$ 推到
$1.0000$。这只表示在冻结的 M1 与成本参数下质量增量成本相对于规模成本较小，
不表示真实训练质量已经饱和，也不表示 $Q_B=1$ 是经验验证的唯一最优质量。

\begin{table}[htbp]
  \centering
  \caption{M1 原生质量条件情景的按预算摘要}
  \label{tab:q3-sp2-summary}
  \small
  \begin{tabular}{lccc}
    \toprule
    质量成本分别为 & $C=10^{19}$ 的平均 $Q_B^*$ & $C=10^{22}$ 的典型状态 & 低预算状态\\
    \midrule
    指数型 & 0.7243 & 接近 $1.0$ & 内点与成本敏感\\
    幂函数型 & 0.6103 & 接近 $1.0$ & 内点与成本敏感\\
    对数渐进型 & 1.0000 & $1.0$ & 质量上界较早活跃\\
    \bottomrule
  \end{tabular}
\end{table}

\subsubsection{上下文成本转移}

将基础成本拆分为训练项 $C_{\mathrm{train}}=6\times10^{18}ND$ 和注意力项
$C_{\mathrm{attn}}=\eta L_{\mathrm{ctx}}\times10^{18}ND$，得到
\begin{equation}
 \frac{C_{\mathrm{attn}}}{C_{\mathrm{train}}}
 =\frac{\eta L_{\mathrm{ctx}}}{6}
 =\frac{L_{\mathrm{ctx}}}{30000}.
 \label{eq:q3-sp2-context-ratio}
\end{equation}
因此 $L_{\mathrm{ctx}}=30000$ 是当前成本结构的代数临界点：低于该值时注意力
项小于训练项，高于该值时注意力项大于训练项。它不是模型效果已经验证的临界点，
也不是质量变量的经验转折点。

\subsubsection{稳健性分析与结论}

复核结果包括：45 行输出及参数折叠传播均为有限值；原生情景最大相对预算误差
约为 $2.1\times10^{-12}$，质量成本稳健性约为 $9.7\times10^{-12}$；
$N,D,Q_B$ 均未越过各自观测支持域；六个起点的收敛和边界状态均被保留；

在 B1 的 $N,D$ 支持域、B6/B7 原生 $Q_B$ 支持域以及给定预算、上下文和成本
假设下，本小问支持报告质量、规模与数据量之间的条件替代关系，以及成本函数和
上下文改变后的资源配置差异。

\subsubsection{小结}

本小问在原生 B6/B7 $Q_B$ 下完成了 45 个条件情景的 $N$--$D$--$Q_B$ 优化。
结果表明，质量成本函数、预算和上下文会共同改变资源配置；参数折叠、多起点、
有限性、预算误差与支持边界检查解除了给定假设下的稳健范围，跨来源
联合最优尚未解决。
